\documentclass[article]{jss}\usepackage[]{graphicx}\usepackage[]{xcolor}
% maxwidth is the original width if it is less than linewidth
% otherwise use linewidth (to make sure the graphics do not exceed the margin)
\makeatletter
\def\maxwidth{ %
  \ifdim\Gin@nat@width>\linewidth
    \linewidth
  \else
    \Gin@nat@width
  \fi
}
\makeatother

\definecolor{fgcolor}{rgb}{0.345, 0.345, 0.345}
\newcommand{\hlnum}[1]{\textcolor[rgb]{0.686,0.059,0.569}{#1}}%
\newcommand{\hlsng}[1]{\textcolor[rgb]{0.192,0.494,0.8}{#1}}%
\newcommand{\hlcom}[1]{\textcolor[rgb]{0.678,0.584,0.686}{\textit{#1}}}%
\newcommand{\hlopt}[1]{\textcolor[rgb]{0,0,0}{#1}}%
\newcommand{\hldef}[1]{\textcolor[rgb]{0.345,0.345,0.345}{#1}}%
\newcommand{\hlkwa}[1]{\textcolor[rgb]{0.161,0.373,0.58}{\textbf{#1}}}%
\newcommand{\hlkwb}[1]{\textcolor[rgb]{0.69,0.353,0.396}{#1}}%
\newcommand{\hlkwc}[1]{\textcolor[rgb]{0.333,0.667,0.333}{#1}}%
\newcommand{\hlkwd}[1]{\textcolor[rgb]{0.737,0.353,0.396}{\textbf{#1}}}%
\let\hlipl\hlkwb

\usepackage{framed}
\makeatletter
\newenvironment{kframe}{%
 \def\at@end@of@kframe{}%
 \ifinner\ifhmode%
  \def\at@end@of@kframe{\end{minipage}}%
  \begin{minipage}{\columnwidth}%
 \fi\fi%
 \def\FrameCommand##1{\hskip\@totalleftmargin \hskip-\fboxsep
 \colorbox{shadecolor}{##1}\hskip-\fboxsep
     % There is no \\@totalrightmargin, so:
     \hskip-\linewidth \hskip-\@totalleftmargin \hskip\columnwidth}%
 \MakeFramed {\advance\hsize-\width
   \@totalleftmargin\z@ \linewidth\hsize
   \@setminipage}}%
 {\par\unskip\endMakeFramed%
 \at@end@of@kframe}
\makeatother

\definecolor{shadecolor}{rgb}{.97, .97, .97}
\definecolor{messagecolor}{rgb}{0, 0, 0}
\definecolor{warningcolor}{rgb}{1, 0, 1}
\definecolor{errorcolor}{rgb}{1, 0, 0}
\newenvironment{knitrout}{}{} % an empty environment to be redefined in TeX

\usepackage{alltt}

%% ---- Packages ----
\usepackage{amsmath,amssymb}
\usepackage{booktabs}
\usepackage{graphicx}
\usepackage{natbib}

%% ---- JSS metadata ----
\author{Arnab Aich\\Florida State University}
\title{DPmixGPD: Bayesian Conditional Dirichlet Process Mixtures with Generalized Pareto Tails and Causal Extensions}

%% Optional: JSS typically wants a short title and keywords.
\Plainauthor{Arnab Aich}
\Plaintitle{DPmixGPD: Bayesian Conditional Dirichlet Process Mixtures with Generalized Pareto Tails and Causal Extensions}
\Shorttitle{DPmixGPD: Conditional DP--GPD Models in R}

\Abstract{
Heavy-tailed outcomes arise in many scientific domains where rare events dominate risk, cost, or policy decisions.
Classical parametric models often fit the center of the distribution at the expense of tail misspecification, while purely nonparametric methods typically lack principled extrapolation in regions with sparse data.
Extreme value theory (EVT) provides asymptotic justification for generalized Pareto distributions (GPDs) for threshold exceedances, whereas Bayesian nonparametric Dirichlet process mixtures (DPMs) provide flexible modeling for heterogeneous bulk behavior.
This article introduces \texttt{DPmixGPD}, an \textsf{R} package implementing a Bayesian semiparametric framework that combines conditional DPMs for the bulk distribution with covariate-dependent GPD tails for exceedances, with posterior predictive inference for densities, distribution functions, quantiles, survival probabilities, and tail risks.
The framework is conditional by construction; unconditional models arise as a special case when covariates are omitted.
The package supports both Chinese restaurant process (CRP) and stick-breaking (SB) representations of the Dirichlet process, multiple kernel families for bulk modeling, and an extension to causal inference via treatment-specific conditional distributions that enable inference on distributional and tail-sensitive causal estimands.
}

\Keywords{Bayesian nonparametrics; Dirichlet process mixture; generalized Pareto distribution; extreme value theory; posterior prediction; causal inference; quantile treatment effects; \textsf{R}; \texttt{nimble}}
\Plainkeywords{Bayesian nonparametrics, Dirichlet process mixture, generalized Pareto distribution, extreme value theory, posterior prediction, causal inference, quantile treatment effects, R, nimble}

%% ---- Begin document ----
\IfFileExists{upquote.sty}{\usepackage{upquote}}{}
\begin{document}
\maketitle

\section{Introduction}
Heavy-tailed outcomes arise in many empirical settings where rare events dominate risk assessment and decision-making.
Examples include environmental and climate applications, where extremes of precipitation or temperature shape planning and infrastructure design \citep{Coles2001,Katz2002,DavisonSmith1990}, insurance and finance where catastrophic losses drive aggregate exposure \citep{Embrechts1997}, and biomedical and public health research where infrequent but severe outcomes may outweigh average effects \citep{AustinSchuster2016}.
Across these domains, analysts frequently require inference for extreme quantiles, tail probabilities, and other functionals of the full outcome distribution rather than only its mean.

Standard parametric models are often inadequate for such tasks because they impose restrictive assumptions on tail behavior.
Misspecification in the tail can induce substantial bias in estimated exceedance probabilities and extreme quantiles even when the model appears adequate in the bulk \citep{Coles2001,DavisonSmith1990}.
Fully nonparametric approaches offer substantial flexibility in regions where data are abundant, but they typically provide limited guidance for extrapolation in the tail where observations are sparse and uncertainty must be carefully controlled \citep{GhosalVanDerVaart2017}.

Extreme value theory (EVT) provides a mathematically principled framework for modeling rare events.
Classical results due to \citet{BalkemaDeHaan1974} and \citet{Pickands1975} show that, under mild regularity conditions, threshold exceedances converge in distribution to a generalized Pareto distribution (GPD).
This asymptotic justification underlies the widespread use of peaks-over-threshold modeling in applied work \citep{DavisonSmith1990,Coles2001,DeHaanFerreira2007}.
In \textsf{R}, EVT methodology is supported by packages such as \texttt{evir} \citep{evirPkg}, \texttt{evd} \citep{Stephenson2002}, \texttt{ismev} \citep{ismevPkg}, and \texttt{extRemes} \citep{GillelandKatz2016}, among others.
However, EVT focuses on extremes and does not specify a model for the bulk distribution, which is often handled heuristically or through ad hoc parametric assumptions.

Bayesian nonparametric methods, particularly Dirichlet process mixtures (DPMs), provide a complementary approach for flexible bulk modeling.
The Dirichlet process introduced by \citet{Ferguson1973} induces a random discrete mixing distribution and has become foundational in Bayesian nonparametrics.
Dirichlet process mixtures are widely used for density estimation and clustering \citep{Lo1984,EscobarWest1995,Neal2000}, and are implemented in \textsf{R} through packages such as \texttt{DPpackage} \citep{Jara2011}, \texttt{dirichletprocess} \citep{RossMarkwick2023}, and \texttt{BNPdensity} \citep{Arbel2021BNPdensity}.
Probabilistic programming tools, notably \texttt{nimble} \citep{DeValpine2017}, further enable custom hierarchical Bayesian models and bespoke MCMC algorithms within \textsf{R}.
Despite these advantages, standard DPMs inherit their tail behavior from the chosen kernel family, which is typically selected for bulk fit rather than for asymptotic tail properties, and this can yield unstable inference for extreme quantiles.

These considerations motivate semiparametric models that combine flexible bulk modeling with EVT-based tails.
Several Bayesian approaches combine nonparametric or flexible bulk models with parametric tails to improve tail inference while retaining bulk flexibility \citep{DoNascimento2012,FuquenePatino2015,CabrasCastellanos2011}.
In particular, \citet{FuquenePatino2015} propose a Dirichlet process mixture of gamma densities for the bulk with a generalized Pareto tail for exceedances, demonstrating stable posterior inference for high quantiles.
However, available implementations are often limited in scope, focus on unconditional models, or do not provide a unified and extensible software suite for conditional modeling, posterior prediction, and causal extensions.

Heavy tails also create challenges for causal inference, where treatment effects may differ across the outcome distribution and where tail behavior can be directly policy-relevant.
Distributional causal estimands such as quantile treatment effects \citep{FirpoFortinLemieux2009,Chernozhukov2013} and tail-risk contrasts can be more informative than average treatment effects when extremes dominate costs or harms.
In \textsf{R}, causal and distributional methods are supported by software such as \texttt{quantreg} \citep{Koenker2005}, \texttt{grf} \citep{AtheyTibshiraniWager2019}, and related packages for propensity score weighting and modern causal modeling.
Nevertheless, many causal workflows rely on outcome models without explicit EVT regularization, which can limit reliability for extreme outcomes.

This article introduces \texttt{DPmixGPD}, an \textsf{R} package implementing a Bayesian semiparametric framework that combines conditional Dirichlet process mixtures for the bulk distribution with covariate-dependent generalized Pareto tails for exceedances, and provides a causal extension via treatment-specific conditional distributions.
The framework is conditional by construction; unconditional models arise as a special case when covariates are omitted.
Posterior inference jointly accounts for uncertainty in the bulk mixture, tail parameters, and the threshold separating bulk and tail regimes, enabling stable posterior predictive inference for densities, distribution functions, quantiles, survival probabilities, and tail risks.

The remainder of the article is organized as follows.
Section~\ref{sec:background} reviews background on conditional DPMs, CRP and stick-breaking representations, and EVT-based tail modeling.
Section~\ref{sec:model} presents the conditional bulk--tail model with a complete likelihood and prior specification.
Section~\ref{sec:posterior} details posterior computation and posterior predictive inference under both CRP and stick-breaking implementations.
Section~\ref{sec:causal} develops the causal extension and the corresponding posterior predictive estimands.
Section~\ref{sec:software} describes the software architecture and user-facing interface of \texttt{DPmixGPD}.
Simulation studies and real-data applications illustrate performance and practical utility.

\section{Background and Related Work}

The methodology implemented in \texttt{DPmixGPD} combines three foundational areas of statistical theory: Bayesian nonparametric mixture modeling for flexible conditional distributions, extreme value theory for tail modeling, and distributional causal inference for principled comparisons across treatment regimes. While each area is well developed individually, they are seldom integrated into a unified inferential and software framework. This section provides a review of the mathematical foundations for each component and describes their integration in \texttt{DPmixGPD}.

\subsection{Conditional Dirichlet Process Mixture Models}

Let $Y \in \mathbb{R}$ denote a univariate outcome, and let $\mathbf{X} = \mathbf{x} \in \mathbb{R}^p$ denote observed covariates. The goal is to estimate the conditional distribution function $F(y \mid \mathbf{x})$ and its associated conditional density $p(y \mid \mathbf{x})$ without restrictive parametric assumptions. Bayesian nonparametric approaches achieve this via flexible mixture models whose complexity adapts to the data.

The conditional Dirichlet process mixture (DPM) model expresses the conditional density as an infinite mixture of kernel densities:
\[
F_{\text{DPM}}(y \mid \mathbf{x}) = \sum_{j=1}^\infty w_j \, F_j(y \mid \mathbf{x}),
\]
where each component distribution $F_j(y \mid \mathbf{x})$ is associated with a kernel $f(y \mid \theta_j(\mathbf{x}))$ and component-specific parameters $\theta_j$, while the weights $\{w_j\}$ arise from a Dirichlet process prior \citep{ferguson1973bayesian, lo1984class}. The covariate dependence enters through $\theta_j(\mathbf{x})$, often modeled using a regression structure with link functions to enforce parameter constraints.

This construction supports localized adaptation of the distributional shape across the covariate space \citep{muller1996bayesian, gelfand2005bayesian}. In \texttt{DPmixGPD}, $F_{\text{DPM}}$ defines the ``bulk'' of the outcome distribution, capturing features such as multimodality and heteroscedasticity in a covariate-dependent manner.

\subsection{CRP and Stick-Breaking Representations}

Two alternative representations of the Dirichlet process are particularly useful for computation: the Chinese restaurant process (CRP) and the stick-breaking (SB) construction. Both are implemented in \texttt{DPmixGPD} to provide algorithmic flexibility.

In the CRP representation, the random mixing measure is integrated out, and inference proceeds by sampling cluster assignments $z_i$ for each observation. The CRP induces the predictive rule:
\[
\Pr(z_i = j \mid z_1,\ldots,z_{i-1}) =
\begin{cases}
\frac{n_j}{i - 1 + \alpha}, & \text{if cluster } j \text{ exists}, \\
\frac{\alpha}{i - 1 + \alpha}, & \text{for a new cluster},
\end{cases}
\]
where $n_j$ is the number of prior observations in cluster $j$. Given $z_i = j$, the likelihood contribution is $f(y_i \mid \theta_j(\mathbf{x}_i))$. This marginal approach often improves mixing in MCMC sampling \citep{neal2000markov}.

Alternatively, the stick-breaking construction represents the random measure explicitly:
\[
G = \sum_{j=1}^\infty w_j \delta_{\theta_j}, \quad
w_j = v_j \prod_{\ell < j} (1 - v_\ell), \quad
v_j \sim \text{Beta}(1, \alpha), \quad
\theta_j \sim G_0,
\]
as introduced by \citet{sethuraman1994constructive}. This form facilitates blocked Gibbs samplers and variational methods \citep{ishwaran2001gibbs}. In \texttt{DPmixGPD}, either representation may be used to define the conditional distribution $F_{\text{DPM}}(y \mid \mathbf{x})$, depending on the computational context.

\subsection{Extreme Value Theory and Generalized Pareto Tails}

While DPM models offer flexibility in the bulk, their tail behavior is determined by the kernel family and is not explicitly modeled. Extreme value theory (EVT) provides a principled framework for modeling the distribution of excesses above a high threshold.

Given a covariate value $\mathbf{x}$, define a threshold $u(\mathbf{x})$, and consider the excess $E = Y - u(\mathbf{x})$ given $Y > u(\mathbf{x})$. EVT shows that, under mild conditions, the conditional distribution of $E$ converges to the generalized Pareto distribution (GPD) \citep{pickands1975statistical,embrechts1997modelling}:
\[
F_{\text{GPD}}(e \mid \sigma(\mathbf{x}), \xi(\mathbf{x})) = 1 - \left(1 + \xi(\mathbf{x}) \frac{e}{\sigma(\mathbf{x})}\right)^{-1/\xi(\mathbf{x})}, \quad e > 0,
\]
with associated density
\[
g(e \mid \sigma(\mathbf{x}), \xi(\mathbf{x})) =
\frac{1}{\sigma(\mathbf{x})}
\left(1 + \xi(\mathbf{x}) \frac{e}{\sigma(\mathbf{x})}\right)^{-1/\xi(\mathbf{x}) - 1},
\]
defined on the support where $1 + \xi(\mathbf{x}) e / \sigma(\mathbf{x}) > 0$. The scale parameter $\sigma(\mathbf{x})$ controls dispersion, while the shape parameter $\xi(\mathbf{x})$ governs tail heaviness.

\texttt{DPmixGPD} models both $\sigma(\mathbf{x})$ and $\xi(\mathbf{x})$ as functions of covariates using link-regression structures, following approaches similar to \citet{coles2001introduction}, allowing tail behavior to vary across the covariate space.

\subsection{Conditional Bulk–Tail Splicing}

To obtain a unified model, \texttt{DPmixGPD} employs a conditional bulk–tail splicing strategy. The overall conditional distribution is defined as:
\[
F(y \mid \mathbf{x}) =
\begin{cases}
F_{\text{DPM}}(y \mid \mathbf{x}), & y \le u(\mathbf{x}), \\
F_{\text{DPM}}(u(\mathbf{x}) \mid \mathbf{x}) +
[1 - F_{\text{DPM}}(u(\mathbf{x}) \mid \mathbf{x})]
F_{\text{GPD}}(y - u(\mathbf{x}) \mid \sigma(\mathbf{x}), \xi(\mathbf{x})), & y > u(\mathbf{x}).
\end{cases}
\]
This formulation guarantees continuity at the threshold and preserves the total probability mass. By integrating nonparametric modeling in the center and parametric modeling in the tail, it enables coherent estimation of the full distribution. Similar splicing constructions have been proposed in both frequentist and Bayesian contexts \citep{tancredi2006bayesian,carreau2009bayesian,naveau2016modeling}.

\subsection{Causal Structure and Modeling Framework}

Causal inference is formulated using the potential outcomes framework \citep{rubin1974estimating,imbens2015causal}. Let $Y^{(1)}$ and $Y^{(0)}$ denote the potential outcomes under treatment and control, respectively, and let $A \in \{0,1\}$ be the observed treatment assignment. We assume a set of baseline covariates $\mathbf{X} = \mathbf{x}$ is available for each unit. Identification of the conditional distributions of potential outcomes relies on the following standard assumptions:

\begin{enumerate}
  \item \textbf{Consistency}: $Y = Y^{(A)}$.
  \item \textbf{No interference}: One unit’s potential outcomes do not depend on the treatment assignment of other units.
  \item \textbf{Conditional exchangeability}: $\{Y^{(0)}, Y^{(1)}\} \perp A \mid \mathbf{X}$.
  \item \textbf{Positivity}: $0 < \Pr(A = 1 \mid \mathbf{X} = \mathbf{x}) < 1$ for all $\mathbf{x}$ in the covariate support.
\end{enumerate}

Under these assumptions, the conditional distributions $F^{(1)}(y \mid \mathbf{x})$ and $F^{(0)}(y \mid \mathbf{x})$ are identifiable from observed data. In \texttt{DPmixGPD}, estimation proceeds by modeling these two potential outcome distributions separately using the same spliced DPM--GPD structure. Covariate adjustment is achieved via propensity score methods \citep{rosenbaum1983central}, where the propensity score $\pi(\mathbf{x}) = \Pr(A = 1 \mid \mathbf{X} = \mathbf{x})$ is estimated and used to adjust for confounding.

This modeling scheme supports flexible, covariate-dependent estimation of both marginal and extreme behavior under each treatment regime, and provides a coherent basis for distributional causal inference.


\section{Model Development: Conditional DP--GPD}
\label{sec:model}
%% TODO: Full likelihood + priors in full sentences.
%% Include both CRP and SB parameterizations as alternative computational representations
%% of the same DP prior.

\section{Posterior Computation and Posterior Predictive Inference}
\label{sec:posterior}
%% TODO: Posterior inference steps + posterior predictive density/CDF/quantile/survival
%% and tail-risk functionals. Include pseudocode + non-executed R code blocks.

\section{Causal Extension}
\label{sec:causal}
%% TODO: Treatment-specific conditional distributions, propensity score augmentation,
%% posterior predictive causal estimands (CQTE, tail contrasts), and computation.

\section{Software}
\label{sec:software}
%% TODO: Package design, main interface functions, supported kernels/backends,
%% prediction API, plots, and reproducibility support.

\section{Simulation Studies}
%% TODO: One-arm conditional simulation and causal simulation with covariates.

\section{Real-Data Applications}
%% TODO: One-arm and causal examples using built-in R datasets (no execution here).

%% ---- Example code chunks (NOT executed) ----
\section*{Appendix: Example Code Skeleton (not executed)}
The following \textsf{R} code chunks illustrate the intended workflow and are included for documentation purposes only.
They are not executed in this manuscript.

\begin{knitrout}
\definecolor{shadecolor}{rgb}{0.969, 0.969, 0.969}\color{fgcolor}\begin{kframe}
\begin{alltt}
\hlcom{# library(DPmixGPD)}

\hlcom{# fit <- dpmixgpd(}
\hlcom{#   y = y,}
\hlcom{#   X = X,}
\hlcom{#   backend = "sb",        # or "crp"}
\hlcom{#   kernel  = "gamma",     # e.g., "normal", "lognormal", "invgauss", "laplace", "cauchy", "amoroso"}
\hlcom{#   gpd     = TRUE,}
\hlcom{#   mcmc    = list(niter = 5000, nburnin = 1000, thin = 5, nchains = 2, seed = c(1,2))}
\hlcom{# )}
\hlcom{# summary(fit)}
\hlcom{# plot(fit)}
\end{alltt}
\end{kframe}
\end{knitrout}

\begin{knitrout}
\definecolor{shadecolor}{rgb}{0.969, 0.969, 0.969}\color{fgcolor}\begin{kframe}
\begin{alltt}
\hlcom{# pred <- predict(}
\hlcom{#   fit,}
\hlcom{#   newdata = Xnew,}
\hlcom{#   type = c("density","cdf","quantile","survival"),}
\hlcom{#   probs = c(0.5, 0.9, 0.99)}
\hlcom{# )}
\hlcom{# pred}
\end{alltt}
\end{kframe}
\end{knitrout}

%% ---- Bibliography ----
\bibliography{DPmixGPD_JSS_refs}

\end{document}
